\section{ORIUS Runtime Architecture}
\label{sec:ieee-runtime}

ORIUS is organized as a post-nominal runtime layer. It does not replace the upstream
predictor, planner, or optimizer. It wraps them with a narrow contract that starts at
raw telemetry and ends at a certificate-gated action release. This keeps the framework
compatible with legacy controllers while moving the safety argument to the place where
degraded observation actually matters.

\begin{figure*}[t]
\centering
\begin{tikzpicture}[
  node distance=0.9cm and 0.7cm,
  stage/.style={draw, rounded corners, minimum width=2.5cm, minimum height=1cm, align=center, fill=gray!8},
  note/.style={font=\small, align=center}
]
\node[stage] (detect) {Detect\\reliability};
\node[stage, right=of detect] (calibrate) {Calibrate\\state set};
\node[stage, right=of calibrate] (constrain) {Constrain\\action set};
\node[stage, right=of constrain] (shield) {Shield\\repair / fallback};
\node[stage, right=of shield] (certify) {Certify\\release};
\draw[-{Latex[length=2.5mm]}] (detect) -- (calibrate);
\draw[-{Latex[length=2.5mm]}] (calibrate) -- (constrain);
\draw[-{Latex[length=2.5mm]}] (constrain) -- (shield);
\draw[-{Latex[length=2.5mm]}] (shield) -- (certify);
\node[note, below=0.65cm of calibrate] {Telemetry, freshness,\\consistency, missingness};
\node[note, below=0.65cm of shield] {Candidate action, repaired action,\\fallback mode, certificate payload};
\end{tikzpicture}
\caption{The ORIUS Detect--Calibrate--Constrain--Shield--Certify kernel. The universal claim is about this runtime grammar rather than about any single nominal controller.}
\label{fig:ieee-runtime-kernel}
\end{figure*}

\subsection{Detect}
The detect stage computes an explicit trust signal from freshness, missingness, timing,
cross-sensor consistency, and domain-specific degradation indicators. The result is not
only an anomaly flag. It is a reliability object that changes the semantics of the rest
of the pipeline.

\subsection{Calibrate}
The calibrate stage inflates the observation-consistent state set. In the current ORIUS
implementation this inflation can be driven by conformal or reliability-conditioned
uncertainty, but the runtime requirement is broader: when observation quality drops, the
set of physically plausible states must not remain frozen at the nominal point estimate.

\subsection{Constrain}
The constrain stage recomputes which actions are admissible under the widened state set.
This is the point at which ORIUS ceases to be an alarming subsystem and becomes a true
safety layer. The system must translate degraded observation into a smaller safe action
surface, not merely record that confidence has fallen.

\subsection{Shield and Certify}
The shield stage repairs the proposed action or falls back to a bounded alternative when
the tightened set becomes empty. The certify stage records why the released action was
allowed, which margins and reliability assumptions were used, and which lifecycle state
the runtime now occupies. That certificate is a scientific object in ORIUS, not a
deployment afterthought.

The architectural value of this decomposition is that the same pipeline appears in every
domain chapter even though the state, action, and repair geometry are different. That
shared pipeline is what makes ORIUS universal in architecture before it is universal in
evidence depth.
